\documentclass[10pt]{article}
\usepackage[margin=0.75in, top=0.65in, bottom=0.65in]{geometry}
\usepackage[compact]{titlesec}
\usepackage{parskip}
\usepackage{listings}
\usepackage{xcolor}
\usepackage{enumitem}
\usepackage[hidelinks]{hyperref}

\setlength{\parskip}{3pt}
\titlespacing{\section}{0pt}{7pt}{3pt}

\lstset{
  basicstyle=\ttfamily\footnotesize,
  backgroundcolor=\color{gray!10},
  frame=single,
  breaklines=true,
  aboveskip=2pt,
  belowskip=2pt
}

\title{\vspace{-1cm}\textbf{CMPT 276: Code Review Report}\vspace{-0.5cm}}
\author{Reviewer: Samantha Gan \quad Reviewee: Krish Kaushik}
\date{}


\begin{document}
\maketitle
\vspace{-1cm}

\section*{Overview}

For Assignment 4, I have reviewed two classes that Krish has contributed to: \texttt{PoliceEnemy.java} and \texttt{GameScreen.java}. I have verified each section of code to be Krish's by reviewing commit history, and ensuring his name is on the blame. Through investigation, I have found 4 code smells for Krish to refactor.


\section{Smell 1: Code Duplication}

\textbf{Files:} \texttt{PoliceEnemy.java} \quad
\textbf{Lines:} \texttt{142-154 and 187-199} \quad
\textbf{Commit:} \texttt{eeac4fb}

Duplicated \texttt{blockedX} and \texttt{blockedY} in alert/chasing state and the wander state. Both compute identical \texttt{hitW, hitTop, blockedX,} and \texttt{blockedY} values using the same four \texttt{isBarrier()} calls each time. 
Lines 142 - 154 and Lines 187 - 199 are the same.

\begin{lstlisting}
boolean blockedX =
            lookup.isBarrier(newX - hitW, position.y - SIZE / 2f) ||
            lookup.isBarrier(newX + hitW, position.y - SIZE / 2f) ||
            lookup.isBarrier(newX - hitW, position.y - hitTop) ||
            lookup.isBarrier(newX + hitW, position.y - hitTop);
        if (!blockedX) position.x = newX;
boolean blockedY = ... if (!blockedY) position.y = newY; // Similar but Y
\end{lstlisting}

This can be fixed by putting the repeated code into a dedicated private method (eg. \texttt{applyBarrierCollision()}).

\section{Smell 2: Method is too long}

\textbf{Files:} \texttt{update()} method in \texttt{PoliceEnemy.java} \quad
\textbf{Lines:} \texttt{96 - 221} \quad
\textbf{Commit:} \texttt{eeac4fb}

The update() method in \texttt{PoliceEnemy.java} handles state, movement/collision and animation. This method is long, and is difficult to read and maintain. 

\begin{lstlisting}
public void update(float delta, float playerX, float playerY, BarrierLookup lookup) {
        // state machine
        // movement ...}
\end{lstlisting}

Having  helper functions for each task can increase readability. By calling each helper in the \texttt{update()} method, it reduces the amount of code inside the method, and breaking it up into 3 digestible, separate chunks. \texttt{update()} serves as a clean coordinator.

Separate into different methods:
\begin{itemize}
    \item \texttt{updateAlertState(float delta, float dist)}
    \item \texttt{updateMovement(float delta, float playerX, float playerY, BarrierLookup lookup)}
    \item \texttt{updateAnimation(float delta)}
\end{itemize} 


\section{Smell 3: Low Cohesion}

\textbf{Files:} \texttt{render()} in \texttt{PoliceEnemy.java} \quad
\textbf{Lines:} \texttt{223-239} \quad
\textbf{Commit:} \texttt{eeac4fb}

The \texttt{render()} method mixes two different responsibilities: sprite drawing, 
and conditional logic that reads \texttt{alertState} to decide which indicator 
to display.

\begin{lstlisting}
public void render(SpriteBatch batch) {
	// sprite drawing
    // draw ? or ! above head depending on state ...}
\end{lstlisting}

A render method should only be presentational, and should not handle any logic (deciding what to display). These two should be separated into \texttt{renderSprite(SpriteBatch batch)} 
and \texttt{renderAlertIndicator(SpriteBatch batch)}, with \texttt{render()} 
calling both. It also makes each piece independently testable and readable.


\section{Smell 4: Data Clump / Magic Numbers}

\textbf{Files:} \texttt{activateMap()} in \texttt{GameScreen.java} \quad
\textbf{Lines:} \texttt{1016-1060} \quad
\textbf{Commit:} \texttt{3c4a96a}

In \texttt{activateMap()}, police and puddle enemies are spawned using long lists 
of hard-coded coordinates with no explanation of what each position represents.

\begin{lstlisting}
police.add(new PoliceEnemy(500f,  mapYOffsets[2] + 850f));
police.add(new PoliceEnemy(1300f, mapYOffsets[2] + 1200f));
// ... 18 more
\end{lstlisting}

There are 20 police spawns and 13 puddle spawns for map 3 alone, all as
float numbers. This is a data clump, where a group of values that always appear together 
but are not encapsulated. If the map layout changes, every coordinate must be 
searched for and updated manually.

The fix is to extract the coordinates into named constant arrays at the top of 
\texttt{GameScreen.java} (e.g. \texttt{MAP3\_POLICE\_SPAWNS}, \texttt{MAP3\_PUDDLE\_SPAWNS}), 
and create a  helper method \texttt{spawnEnemiesForMap(int index)} that 
iterates over the arrays. \texttt{activateMap()} then calls 
\texttt{spawnEnemiesForMap(index)}, making it shorter and easier to maintain.

\begin{lstlisting}
private static final float[][] MAP3_POLICE_SPAWNS = {
    {500f,  850f}, {1300f, 1200f}, {2500f, 1300f}, ...
};
\end{lstlisting}


\section*{Refactoring 
(Based on Megan's Review)}

After receiving Megan's code review, I refactored each identified code smell as follows:

\textbf{Smell 1: Unnecessary Variable} \quad \textbf{Commit:} \texttt{fd23f22} \\
\texttt{bikeParked} was removed from \texttt{GameScreen.java}. All assignments 
to \texttt{bikeParked} were deleted, and the one usage \texttt{if (bikeParked)} 
in the render loop was replaced with \texttt{if (!isOnBike)}. The game was 
rerun and bike parking, mounting, and dismounting all behaved correctly.

\textbf{Smell 2: Data Clump} \quad \textbf{Commit:} \texttt{cc410b2} \\
In \texttt{EndScreen.java}, the \texttt{DIGIT\_} and \texttt{TIME\_} constants 
were grouped into a \texttt{RenderLayout} structure so that related values 
(position, size, gap) are encapsulated together rather than stored as separate 
unrelated fields. The end screen was tested and displayed correctly after the change.

\textbf{Smell 3: Code Duplication} \quad \textbf{Commit:} \texttt{7208b9f} \\
In \texttt{EndScreen.java}, the repeated digit-extraction pattern (\texttt{/10}, 
\texttt{\%10}) shared by \texttt{scoreDigits} and \texttt{timeDigits} was 
extracted into a private helper method. Both arrays now call the shared method 
instead of repeating the arithmetic. The end screen was tested and displayed 
correct score and time values.

\textbf{Smell 4: Dead Code} \quad \textbf{Commit:} \texttt{9782963} \\
In \texttt{House.java}, the unused import was removed. No behaviour changed 
as the import was never referenced. The project was rebuilt successfully.

\end{document}
